%%%%%%%%%%%%%%%%%%%%%%%%%%%%%%%%%%%%%%%%%%%%%%%%%%%%%%%%%%%%%%%%%%%%
\section{Proposed Method}
\label{sec:method}

\paragraph{Overview.}
\Cref{fig:arch} gives the full pipeline. A frozen OVD with a generic prompt first finds product boxes; it makes no attempt to name the SKU at this point. Each crop then goes through a frozen DINOv2 encoder. On the catalogue side, every SKU has one to five reference images, embedded with the same encoder, and a name, embedded with a frozen CLIP text encoder. An adapter pulls near-identical SKUs apart in feature space, and a re-ranker merges visual, text and detector evidence into one score per SKU. Together they hold about 0.7\,M parameters, so a new SKU needs a reference photo and nothing else.

\paragraph{(a) Class-agnostic proposals.}
The input is a shelf image $I$. We give the frozen Grounding DINO \cite{liu2024groundingdino} the prompt ``product . box . pouch . bottle .'' and keep the boxes $b_i$, $i=1,\dots,N$, whose confidence $s^{\mathrm{det}}_i$ exceeds 0.25 after non-maximum suppression at IoU 0.5. To see how much rides on that choice, we also run YOLO-World \cite{cheng2024yoloworld} as the proposer. Nothing in this stage is trained. Proposal recall is reported on its own, so a box the detector missed is never counted as a wrong SKU.

\paragraph{(b) Frozen embeddings.}
We crop each box with a 10\% margin, resize it to $224\times224$ and pass it through DINOv2 ViT-B/14 \cite{oquab2024dinov2}, taking the class token as the crop feature $f_i\in\mathbb{R}^{768}$. The reference images of SKU $k$ go through the same encoder and their features are averaged into $e_k$. We picked DINOv2 over CLIP here because self-supervised features tend to keep the small print and texture that caption-based training washes out, and \cref{sec:plan} checks that assumption. The text branch uses CLIP ViT-B/16 \cite{radford2021clip}. The SKU name, written as ``brand, line, variant, size'', gives $t_k$, and the crop passed through the CLIP image encoder gives $c_i$, so $\cos(c_i,t_k)$ acts as a zero-shot text match. None of the encoders are updated.

\paragraph{(c1) SKU adapter (trained).}
The adapter is a two-layer residual MLP, $g_\theta(x)=W_2\,\mathrm{GELU}(W_1x)+Px$, with $W_1\in\mathbb{R}^{512\times768}$, $W_2\in\mathbb{R}^{256\times512}$ and a linear shortcut $P\in\mathbb{R}^{256\times768}$, followed by $\ell_2$ normalisation. Crops become $z_i=g_\theta(f_i)$ and references become $\hat z_k=g_\theta(e_k)$. Using one adapter for both is deliberate. A SKU that never appeared in training can still be matched at test time, because its reference photo is mapped by the same function as the crops.

\paragraph{(c2) Fusion re-ranker (trained).}
For crop $i$ and SKU $k$ the score is
\begin{equation}
  s_{ik} = w_v \cos(z_i,\hat z_k) + w_t \cos(c_i,t_k) + w_d\, s^{\mathrm{det}}_i ,
  \label{eq:score}
\end{equation}
with three learned scalars. The predicted SKU is $\arg\max_k s_{ik}$; if $\max_k s_{ik}<\delta$, the crop is labelled \emph{unknown} (panel d of \cref{fig:arch}). We set $\delta$ on the validation split to give 95\% precision on known SKUs. Some reject rule is needed, since every real shelf carries products that are not in the catalogue.

\paragraph{(c3) Objective.}
The adapter and the fusion weights are trained together with a contrastive loss over the gallery $\mathcal{K}$ of base SKUs. SKUs that share a brand with the true SKU $y_i$ get extra weight in the denominator:
\begin{equation}
  \mathcal{L} = -\frac{1}{N}\sum_{i=1}^{N}\log\frac{\exp(s_{iy_i}/\tau)}{\sum_{k\in\mathcal{K}}\omega_{ik}\exp(s_{ik}/\tau)} ,
  \label{eq:loss}
\end{equation}
where $\omega_{ik}=\gamma$ if $k\neq y_i$ has the brand of $y_i$ and $\omega_{ik}=1$ otherwise, and $\tau$ is a learned temperature. With $\gamma=1$ this is the standard InfoNCE loss \cite{oord2018cpc}. Raising $\gamma$ pushes the gradient towards the near-duplicates that generic features confuse, an idea close to supervised contrastive learning \cite{khosla2020supcon}. Training crops are cut from ground-truth boxes in the training split and augmented with random resized crops, colour jitter and perspective warps, so the adapter sees something like the gap between a studio photo and a shelf.
